\documentclass[11pt,letterpaper]{article}
\usepackage[T1]{fontenc}
\usepackage[utf8]{inputenc}
\usepackage{lmodern}
\usepackage[margin=1in,headheight=14pt]{geometry}
\usepackage{microtype}
\usepackage{amsmath,amssymb,amsthm,mathtools}
\usepackage{booktabs,array,longtable}
\usepackage{aliascnt}
\usepackage{xcolor}
\usepackage{enumitem}
\usepackage{fancyhdr}
\usepackage[most]{tcolorbox}
\usepackage{listings}
\usepackage{hyperref}
\usepackage[nameinlink,noabbrev]{cleveref}
\definecolor{Navy}{HTML}{17344D}
\definecolor{Teal}{HTML}{157F82}
\definecolor{Pale}{HTML}{EFF6F7}
\definecolor{Gray}{HTML}{53606B}
\hypersetup{colorlinks=true,linkcolor=Navy,citecolor=Teal,urlcolor=Teal,
 pdftitle={Open-query games on ordinal spaces: a negative solution to the topological-sum problem},
 pdfauthor={Research report},
 pdfsubject={Order topology, Cantor--Bendixson rank, finite decision trees}}
\pagestyle{fancy}
\fancyhf{}
\fancyhead[L]{\small\sffamily\color{Navy} ORDINAL MEMBERSHIP GAMES}
\fancyhead[R]{\small\sffamily\color{Gray} RESEARCH REPORT}
\fancyfoot[C]{\small\thepage}
\renewcommand{\headrulewidth}{0.3pt}
\setlength{\parskip}{0.35em}
\setlength{\parindent}{1.2em}
\setlist{nosep,leftmargin=1.7em}
\numberwithin{equation}{section}
\newtheorem{theorem}{Theorem}[section]
\newaliascnt{proposition}{theorem}
\newtheorem{proposition}[proposition]{Proposition}
\aliascntresetthe{proposition}
\newaliascnt{lemma}{theorem}
\newtheorem{lemma}[lemma]{Lemma}
\aliascntresetthe{lemma}
\newaliascnt{corollary}{theorem}
\newtheorem{corollary}[corollary]{Corollary}
\aliascntresetthe{corollary}
\theoremstyle{definition}
\newaliascnt{definition}{theorem}
\newtheorem{definition}[definition]{Definition}
\aliascntresetthe{definition}
\newaliascnt{example}{theorem}
\newtheorem{example}[example]{Example}
\aliascntresetthe{example}
\theoremstyle{remark}
\newaliascnt{remark}{theorem}
\newtheorem{remark}[remark]{Remark}
\aliascntresetthe{remark}
\newcommand{\sm}{\operatorname{sm}}
\newcommand{\ps}{\operatorname{ps}}
\newcommand{\cl}{\operatorname{cl}}
\newcommand{\rk}{\operatorname{rk}}
\newcommand{\ceil}[1]{\left\lceil#1\right\rceil}
\newcommand{\N}{\mathbb N}
\newcommand{\eps}{\varepsilon}
\newcommand{\bigsqsum}{\bigsqcup}
\newcommand{\Fin}{\operatorname{Fin}}
\newcommand{\CB}{\operatorname{CB}}
\newtcolorbox{resultbox}{colback=Pale,colframe=Teal,boxrule=0.5pt,
 arc=1mm,left=3mm,right=3mm,top=2mm,bottom=2mm,breakable}
\lstset{basicstyle=\small\ttfamily,columns=fullflexible,keepspaces=true,
 breaklines=true,frame=single,rulecolor=\color{Gray},showstringspaces=false}

\begin{document}
\hypersetup{pageanchor=false}
\begin{titlepage}
\vspace*{0.4in}
{\sffamily\color{Teal}\large ORDER TOPOLOGY \; / \; ORDINAL COMBINATORICS\par}
\vspace{0.45in}
{\sffamily\bfseries\color{Navy}\Huge Open-query games\\[0.15em] on ordinal spaces\par}
\vspace{0.3in}
{\sffamily\color{Navy}\Large A negative solution to the topological-sum problem,\\
with an exact Cantor--Bendixson classification\par}
\vspace{0.5in}
{\large Research report\par}
{\color{Gray}19 September 2026\par}
\vspace{0.45in}
\begin{resultbox}
For the convergent ordinal sequence $S=[0,\omega]$,
\[
 \sm(S)=1,
 \qquad
 \sm(S\sqcup S)=2.
\]
More generally, for every integer $k\geq1$, the compact ordinal space
\[
 K_k=[0,\omega^{\,2^k-1}]
\]
satisfies $\sm(K_k)=k$ and $\sm(K_k\sqcup K_k)=k+1$.
\end{resultbox}
\vfill
{\small\color{Gray}
Complete ordinary mathematical proofs are included. The archive also contains
an independent finite-poset verifier, exact data, and a machine-readable strategy
certificate. The infinite-space results do not depend on numerical experiments.
No claim of independent peer review or formal proof-assistant verification is made.\par}
\end{titlepage}
\hypersetup{pageanchor=true}

\begin{abstract}
We answer negatively Problem 1.3 of Chiozini, Csern\'ak, and Soukup's
\emph{Gamification of the $T_0$-pseudoweight via cut-and-choose games on
topological spaces}, in the version available as arXiv:2510.05754v3.
The counterexample consists of two convergent sequences, so it lies entirely
within compact metrizable ordinal spaces.

The main further result is an exact formula for every nonempty Hausdorff space
of finite Cantor--Bendixson height. If $X^{(h)}\ne\varnothing$ and
$X^{(h+1)}=\varnothing$, then
\[
 \sm(X)=
 \begin{cases}
  \ceil{\log_2(h+1)},& |X^{(h)}|=1,\\
  \ceil{\log_2(h+2)},& |X^{(h)}|\geq2.
 \end{cases}
\]
The proof identifies finite open-query complexity with the logarithm of the
minimum length of an alternating closed-layer decomposition. A canonical
alternating-closure procedure computes that length and retains the orientation
information lost by the scalar invariant $\sm$.

Consequences include sharp one-move failures of the proposed sum identity at
every positive finite level, exact finite-product formulae, a classification on
countable scattered Hausdorff spaces, and a quadratic-time target algorithm for
finite posets. Independent minimax computations verify the finite normal form
on all naturally labelled posets through six elements and all their targets:
320,866 tests in total.
\end{abstract}

\setcounter{tocdepth}{1}
{\small\setlength{\parskip}{0pt}\tableofcontents}
\newpage

\section{The selected problem and the status of the answer}

In the set-membership game, questions are restricted to open subsets of a space,
and the aim is to decide membership in a fixed target subset. The invariant
$\sm(X)$ is the worst, over all targets, of the least ordinal number of questions
needed. Chiozini, Csern\'ak, and Soukup ask whether
\begin{equation}\label{eq:question}
 \sm\!\left(\bigsqsum_{i\in I}X_i\right)
 \stackrel{?}{=}\sup_{i\in I}\sm(X_i).
\end{equation}
This is their Problem 1.3; their Theorem 3.4 gives the upper bound consisting of
the right-hand side followed by one additional move \cite{CCS}.

The answer to \eqref{eq:question} is \emph{no}. The proof in the next section is
short and does not depend on any later theorem. In particular, the counterexample
is not an artifact of non-Hausdorff topology, uncountable spaces, or an ambiguity
about transfinite play. It uses a two-move game on a countable compact metrizable
space. Repeating the same component countably many times also gives a
counterexample when the index set is required to be infinite.

The more substantial development is the classification in
\cref{thm:height}. Its mechanism explains exactly why the counterexample works.
Different components may require opposite interpretations of the same
open-set answer. A scalar question count forgets that orientation, while a pair
of alternating-closure lengths preserves it.

\paragraph{Status and provenance.}
The arXiv record retrieved for this report lists version 3, dated 29 May 2026;
the question remains explicitly printed there \cite{CCS}. The report gives a
negative solution to that stated question, not a solution to every open problem
in the paper. A targeted literature search found no intervening resolution of
this question. This is a search statement, not a proof of priority. The
closed-layer machinery belongs to the classical setting of Hausdorff difference
hierarchies; we do not claim that framework as new \cite{BGKS}. All lemmas needed
for the results here are proved below.

\section{A compact ordinal counterexample}\label{sec:small}

Let $S=[0,\omega]$ have its order topology. Thus every natural number is isolated,
and the neighborhoods of $\omega$ contain a final segment of $\N$.
We use zero questions for a constant target: a subset equal to the whole space
or to the empty set.

\begin{lemma}[The one-question criterion]\label{lem:one}
For a subset $A$ of any topological space $X$, membership in $A$ can be decided
with at most one open-set question if and only if $A$ is open or closed.
\end{lemma}
\begin{proof}
An open target is itself a legal question. A closed target can be tested by
asking its open complement. Conversely, suppose the first question is $U$.
The only two possible final states are $U$ and $X\setminus U$. Both must be
monochromatic for membership in $A$. If they have the same color, $A$ is
constant. Otherwise $A=U$ or $A=X\setminus U$.
\end{proof}

\begin{proposition}\label{prop:small}
For $S=[0,\omega]$,
\begin{equation}\label{eq:small}
 \sm(S)=1,\qquad \sm(S\sqcup S)=2.
\end{equation}
\end{proposition}
\begin{proof}
Every subset of $S$ missing $\omega$ is open, since it consists entirely of
isolated points. Every subset containing $\omega$ is closed, since its complement
consists of isolated points. Hence every target takes at most one question.
There are nonconstant targets, so $\sm(S)=1$.

Write $S_0,S_1$ for two disjoint copies, with limit points $p_0,p_1$, and put
\[
 A=\{p_0\}\ \sqcup\ (S_1\setminus\{p_1\}).
\]
This set is not open: it contains $p_0$ but no isolated point of $S_0$.
It is not closed: it contains all isolated points of $S_1$, but omits their
limit $p_1$. By \cref{lem:one}, one question cannot suffice.

Two questions suffice for \emph{every} target $B$. First ask the clopen set
$S_0$, thereby determining which copy contains the remaining candidates.
Inside that copy, $B$ is open or closed, so one more open question decides
membership. An open set in a component is open in the disjoint union.
\end{proof}

The fact that spaces with at most one accumulation point have one-question
membership complexity is also noted in \cite[Section 7]{CCS}; the failure of
\eqref{eq:question} comes from putting two such components together.

\begin{corollary}
The proposed identity fails even when every component and the whole sum are
compact, countable, metrizable, zero-dimensional, and homeomorphic to ordinal
spaces. Specifically,
\[
 S\sqcup S\cong[0,\omega\cdot2],
 \qquad \sm([0,\omega\cdot2])=2.
\]
\end{corollary}
\begin{proof}
The two clopen intervals $[0,\omega]$ and $(\omega,\omega\cdot2]$ partition
$[0,\omega\cdot2]$. Each has order type $\omega+1$. The stated topological
properties can also be seen directly from two disjoint convergent sequences.
\end{proof}

\begin{proposition}[An infinite index set]
For every cardinal $\kappa\geq2$,
\[
 \sm\!\left(\bigsqsum_{i<\kappa}S\right)=2.
\]
\end{proposition}
\begin{proof}
The lower bound follows by restricting to two components. For the upper bound,
fix a target $A$, and split the component indices into two classes so that $A_i$
is open on the first class and closed on the second; when both apply, choose
either class. One question asks the union of the components in the first class.
A second question asks the open set whose trace is $A_i$ on the first class and
$S_i\setminus A_i$ on the second. The two answers determine the color.
\end{proof}

\section{The game, subspaces, and ordinal conventions}

All ordinal intervals in this report carry their order topology. The notation
$[0,\gamma]$ denotes a space whose underlying ordered set has order type
$\gamma+1$. A topological sum $X\sqcup Y$ is not being confused with an ordinal
sum; when they are homeomorphic, this is stated separately.

\begin{definition}[Open-query membership game]
Fix $A\subseteq X$ and an ordinal $\beta$. At stage $\xi<\beta$, the Seeker
chooses an open set $U_\xi$. The Hider answers $e_\xi\in\{0,1\}$. Set
\[
 W_\eta=\bigcap_{\xi<\eta}
 \begin{cases}
  U_\xi,&e_\xi=1,\\
  X\setminus U_\xi,&e_\xi=0.
 \end{cases}
\]
The Seeker wins if $W_\beta\subseteq A$ or $W_\beta\cap A=\varnothing$.
Write $\sm(A,X)$ for the least winning length and
$\sm(X)=\sup_{A\subseteq X}\sm(A,X)$, as in \cite{CCS}.
\end{definition}

For ordinal-valued invariants we work with $T_0$ spaces, where a winning
length exists: asking all open sets separates points. A statement of the form
$\sm(A,X)\leq d$ can also be read on an arbitrary topological space as the
existence of a winning depth-$d$ strategy, whether or not a least ordinal
winning length exists. The finite normal form is valid in this broader sense.

An empty final state is winning. The Hider need not commit to a point in advance.
For a finite decision tree, however, the game formulation and correctness on
every fixed point are equivalent: a bad leaf contains two points of different
colors, while an empty leaf is harmless. This equivalence will be used carefully
in \cref{sec:normal}.

\begin{lemma}[Restriction to a subspace]\label{lem:subspace}
If $Z\subseteq X$ and $A\subseteq X$, then
\[
 \sm(A\cap Z,Z)\leq\sm(A,X),\qquad \sm(Z)\leq\sm(X).
\]
\end{lemma}
\begin{proof}
Restrict every open question in a strategy on $X$ to $Z$. It is open in the
subspace, and every resulting state is the old state intersected with $Z$.
For the second inequality, extend an arbitrary target of $Z$ to a target of $X$.
\end{proof}

Finite question counts are ordinary nonnegative integers. We use
$\ceil{\log_2 L}$ only for positive integers $L$; equivalently, it is the least
integer $d$ for which $L\leq2^d$. This avoids floating-point or ordinal-logarithm
conventions. In particular, $\ceil{\log_2 1}=0$.

\section{Finite decision trees and alternating layers}\label{sec:normal}

\begin{definition}[Alternating closed layers]
For a nonempty space $X$ and $A\subseteq X$, let $C_1=A$ and
$C_0=X\setminus A$. A closed-layer decomposition of length $L\geq1$ and
outer color $c\in\{0,1\}$ is a chain
\begin{equation}\label{eq:layerchain}
 X=F_0\supseteq F_1\supseteq\cdots\supseteq F_L=\varnothing
\end{equation}
of closed sets such that
\begin{equation}\label{eq:layercolor}
 F_i\setminus F_{i+1}\subseteq C_{(c+i)\bmod2}
 \quad(0\leq i<L).
\end{equation}
Empty layers are allowed. Let $L(A,X)$ be the least such length, allowing either
outer color, and put $L(A,X)=\infty$ if no finite chain exists.
\end{definition}

These are finite alternating differences of closed sets, equivalently finite
Boolean combinations of open sets. This connects the game to the difference
hierarchy, but the exact decision-depth statement below is proved directly.

\begin{theorem}[Finite normal form]\label{thm:normal}
For every nonempty topological space $X$, target $A$, and integer $d\geq0$,
\begin{equation}\label{eq:normal}
 \sm(A,X)\leq d
 \quad\Longleftrightarrow\quad L(A,X)\leq2^d.
\end{equation}
Consequently, whenever $L(A,X)<\infty$,
\begin{equation}\label{eq:logdepth}
 \sm(A,X)=\ceil{\log_2 L(A,X)}.
\end{equation}
\end{theorem}

\subsection{From a query tree to an ordered collection of leaves}

\begin{lemma}\label{lem:leaf}
Complete a finite open-query tree to depth $d$, padding terminal branches with
dummy questions. Number its leaves $0,1,\ldots,2^d-1$ in lexicographic order of
answer strings, with $0<1$. If $\ell(x)$ is the leaf reached by the answers of
$x$, then every upper level set
\[
 \{x:\ell(x)\geq j\}
\]
over the leaf order is open.
\end{lemma}
\begin{proof}
Proceed by induction on $d$. The case $d=0$ is immediate. Let the root question
be $U$, and put $M=2^{d-1}$. Consider the two subtrees as query trees defined on
all of $X$, regardless of which root answer would actually occur, and write
$\ell_0,\ell_1$ for their leaf functions. If $j\leq M$, then
\[
 \{\ell\geq j\}=U\cup\{\ell_0\geq j\};
\]
for $j=M$, the second set is empty. If $j>M$, then
\[
 \{\ell\geq j\}=U\cap\{\ell_1\geq j-M\}.
\]
The induction hypothesis and closure of the topology under finite unions and
intersections prove the claim.
\end{proof}

\begin{proof}[Proof of the forward implication of \cref{thm:normal}]
A winning depth-$d$ tree has a well-defined output color at each nonempty leaf.
Assign an arbitrary color to each empty leaf. Read the $2^d$ leaf colors from
the largest leaf down to the smallest, and group consecutive equal colors into
runs. There are at most $2^d$ runs, and successive run colors alternate.

After deleting the largest $i$ runs, the set of points still represented is a
lower level set of $\ell$. It is closed by \cref{lem:leaf}. These lower level
sets, starting with $X$ and ending with $\varnothing$, form the required chain.
A run can represent no points, which is why empty layers were allowed.
\end{proof}

\subsection{Balanced search among closed layers}

\begin{proof}[Proof of the reverse implication of \cref{thm:normal}]
Suppose \eqref{eq:layerchain}--\eqref{eq:layercolor} hold. Every point lies in a
unique layer $F_i\setminus F_{i+1}$. The open question $X\setminus F_j$ tests
whether $i<j$. Thus ordinary binary search can identify the layer index among
$L$ possibilities. Split the current interval of indices as evenly as possible
at each question. Its largest remaining size after $d$ questions is at most
$\ceil{L/2^d}$. When $L\leq2^d$, at most one layer remains, and that layer is
monochromatic. The same reasoning applies to arbitrary Hider answers: either
the state is empty or it is contained in a single layer.
\end{proof}

\begin{remark}[Why the logarithm is genuine]
Sequentially peeling $L$ layers would spend roughly $L$ questions. Adaptivity
permits balanced search through nested open thresholds, reducing the count to
$\ceil{\log_2 L}$. Conversely, a depth-$d$ tree has at most $2^d$ ordered leaf
runs. Neither direction assumes that the space is ordered or Hausdorff.
\end{remark}

\section{Canonical alternating closures and orientation}

The least layer length has a concrete canonical description. This is also the
right quantity for understanding disjoint unions.

\begin{definition}[Oriented closure sequence]\label{def:closure}
For $c\in\{0,1\}$ define
\begin{equation}\label{eq:greedy}
 F^c_0=X,\qquad
 F^c_{i+1}=\cl_X\bigl(F^c_i\cap C_{(c+i+1)\bmod2}\bigr).
\end{equation}
Let $N_c(A,X)$ be the first integer $L$ such that $F^c_L=\varnothing$, or
$\infty$ when no such integer exists. The pair $(N_0,N_1)$ is the
\emph{oriented closure profile} of the target.
\end{definition}

Every $F^c_i$ is closed and the sequence is decreasing. At step $i+1$, one keeps
the closure of points whose color would be wrong in the next layer. Therefore
\[
 F^c_i\setminus F^c_{i+1}\subseteq C_{(c+i)\bmod2}.
\]

\begin{theorem}[Canonical minimality]\label{thm:greedy}
For each outer color $c$, $N_c(A,X)$ is the least possible length of a
closed-layer decomposition with that color. In particular,
\begin{equation}\label{eq:profileformula}
 L(A,X)=\min(N_0(A,X),N_1(A,X)).
\end{equation}
\end{theorem}
\begin{proof}
If the canonical sequence terminates, its layers are a decomposition of the
required length. Conversely, let $G_0,\ldots,G_L$ be any closed-layer chain
with outer color $c$. We show $F^c_i\subseteq G_i$ by induction.
The assertion is true at $i=0$. If it holds at $i$, then every wrong-colored
point of $F^c_i$ belongs to $G_{i+1}$, because $G_i\setminus G_{i+1}$ has the
prescribed color. Taking closures preserves containment in the closed set
$G_{i+1}$. Hence $F^c_{i+1}\subseteq G_{i+1}$. In particular $F^c_L=\varnothing$.
\end{proof}

\begin{lemma}[Two elementary profile rules]\label{lem:profiles}
Complementing the target swaps $N_0$ and $N_1$. Also,
\[
 N_c\leq N_{1-c}+1
\]
whenever the latter is finite.
\end{lemma}
\begin{proof}
The first assertion follows directly from \eqref{eq:greedy}. For the second,
prepend an empty outer layer of color $c$ to a decomposition whose outer color
is $1-c$; in the chain this means repeating $X$ once at the beginning.
Apply \cref{thm:greedy}.
\end{proof}

Thus finite profile entries differ by at most one, but discarding the distinction
between them can change a question count at a power-of-two threshold.

\begin{proposition}[Exact finite profile of a sum]\label{prop:profile-sum}
Let $X=\bigsqsum_{i\in I}X_i$, $A_i=A\cap X_i$, and let $c\in\{0,1\}$.
For a nonempty family of nonempty components,
\begin{equation}\label{eq:profile-sum}
 N_c(A,X)=\sup_{i\in I}N_c(A_i,X_i),
\end{equation}
where the supremum is taken in $\N\cup\{\infty\}$: an unbounded set of finite
lengths has value $\infty$. Hence, whenever the resulting minimum is finite,
\begin{equation}\label{eq:sumexact}
 \sm(A,X)=\ceil{\log_2\!\left(
   \min_{c\in\{0,1\}}\sup_i N_c(A_i,X_i)\right)}.
\end{equation}
\end{proposition}
\begin{proof}
Closure in a disjoint union is computed componentwise. Inductively,
$F^c_n(A,X)\cap X_i=F^c_n(A_i,X_i)$ for every finite $n$. The global sequence
is empty exactly when every component sequence is empty. Now use
\cref{thm:normal,thm:greedy}.
\end{proof}

The obstruction to \eqref{eq:question} is already visible in the elementary
inequality
\[
 \min_c\sup_i N_{i,c}\ \geq\ \sup_i\min_c N_{i,c}.
\]
The two operations need not commute: each component may have its own preferred
outer color, whereas one global closed-layer chain has only one outer color.

\section{Exact classification at finite Cantor--Bendixson height}
\label{sec:height}

For a closed subspace $F$ of a Hausdorff space, let $F'$ be its set of
nonisolated points. Set $X^{(0)}=X$, $X^{(n+1)}=(X^{(n)})'$; at a limit ordinal,
use the intersection of the previous derivatives. Only finite derivatives are
needed for the principal formula.

\begin{definition}
A nonempty Hausdorff space has \emph{last derivative rank $h$} if
\[
 X^{(h)}\ne\varnothing,\qquad X^{(h+1)}=\varnothing.
\]
Its Cantor--Bendixson height is then $h+1$. The terminal derivative $X^{(h)}$
is discrete; it need not be finite unless the space is compact.
\end{definition}

\begin{theorem}[Finite-height classification]\label{thm:height}
Let $X$ be a nonempty Hausdorff space with last derivative rank $h<\omega$.
Then
\begin{equation}\label{eq:height}
 \boxed{\quad
 \sm(X)=
 \begin{cases}
  \ceil{\log_2(h+1)},& |X^{(h)}|=1,\\[2pt]
  \ceil{\log_2(h+2)},& |X^{(h)}|\geq2.
 \end{cases}\quad}
\end{equation}
More precisely, the maximum of $L(A,X)$ over all targets is $h+1$ in the first
case and $h+2$ in the second. The maxima are attained.
\end{theorem}

\subsection{Why canonical closures lose one derivative at each step}

\begin{lemma}\label{lem:isolation}
If $F=\cl_X B$ and $x$ is isolated in $F$, then $x\in B$.
\end{lemma}
\begin{proof}
An open neighborhood $U$ with $U\cap F=\{x\}$ must meet $B$, because
$x\in\cl_X B$. Thus $x\in B$.
\end{proof}

\begin{lemma}[Derivative descent]\label{lem:descent}
For every target and outer color, the canonical sequence satisfies
\begin{equation}\label{eq:descent}
 F^c_{n+1}\subseteq(F^c_n)'\quad(n\geq1),
 \qquad F^c_n\subseteq X^{(n-1)}\quad(n\geq1).
\end{equation}
\end{lemma}
\begin{proof}
By \cref{lem:isolation}, every isolated point of $F^c_n$ has color
$(c+n)\bmod2$. Therefore the opposite-colored points whose closure defines
$F^c_{n+1}$ all lie in $(F^c_n)'$. This derivative is closed in $X$, so their
closure lies there too. Derivation is monotone on closed sets, giving the
second assertion by induction from $F^c_1\subseteq X$.
\end{proof}

\begin{proof}[Upper bounds in \cref{thm:height}]
For every target and orientation, \cref{lem:descent} gives
$F^c_{h+2}\subseteq X^{(h+1)}=\varnothing$. Hence $L(A,X)\leq h+2$.

Now assume $X^{(h)}=\{p\}$, and write $d$ for the target color of $p$.
Choose $c$ so that $(c+h)\bmod2=d$. We have $F^c_{h+1}\subseteq\{p\}$.
If this set were nonempty, $p$ would be isolated in it and, by
\cref{lem:isolation}, would have color $(c+h+1)\bmod2=1-d$, a contradiction.
Thus $F^c_{h+1}=\varnothing$ and $L(A,X)\leq h+1$.
\end{proof}

\subsection{Alternating rank layers give matching lower bounds}

Write
\[
 D_j=X^{(j)}\setminus X^{(j+1)}\qquad(0\leq j\leq h).
\]
These sets partition $X$. The set $D_j$ is dense in $X^{(j)}$. Here are details
that do not assume compactness: a finite-height space is scattered, because any
nonempty subspace contains a point of least ambient rank, and that point is
isolated in the subspace after restricting an isolating neighborhood at its rank.
Every nonempty open subset of a scattered space contains a point isolated in
that space. Apply this to $X^{(j)}$.

\begin{lemma}[Parity target]\label{lem:parity}
Color every point of $D_j$ by $(b+j)\bmod2$, where $b\in\{0,1\}$, and let
$A_b$ be the color-$1$ set. Then
\begin{equation}\label{eq:parityprofiles}
 N_b(A_b,X)=h+1,\qquad N_{1-b}(A_b,X)=h+2.
\end{equation}
\end{lemma}
\begin{proof}
For $h=0$, the target is constant of color $b$, and the two lengths are
$1$ and $2$, as asserted. Assume $h\geq1$.
With outer color $b$, the first closure is $X'$, because the opposite color
contains $D_1$, which is dense in $X'$, and is contained in $X'$.
Repeating the argument shows $F^b_n=X^{(n)}$ for $0\leq n\leq h+1$.
Thus the first empty term has index $h+1$.

With outer color $1-b$, the first closure is $X$, since it includes the dense
rank-zero layer $D_0$. Subsequent closures give
$F^{1-b}_n=X^{(n-1)}$ for $1\leq n\leq h+2$. The first empty term now has
index $h+2$. The same computation covers $h=0$.
\end{proof}

\begin{lemma}[Restriction of profiles]\label{lem:profile-restrict}
For any subspace $Z\subseteq X$ and either color $c$,
\[
 N_c(A\cap Z,Z)\leq N_c(A,X).
\]
\end{lemma}
\begin{proof}
The closure identity
$\cl_Z(B\cap Z)\subseteq Z\cap\cl_X B$ gives, inductively,
$F^c_n(A\cap Z,Z)\subseteq Z\cap F^c_n(A,X)$.
\end{proof}

\begin{proof}[Lower bounds in \cref{thm:height}]
The parity target of \cref{lem:parity} always has $L=h+1$, proving the required
lower bound when the terminal derivative is a singleton.

Suppose instead that $p\ne q$ lie in $X^{(h)}$. Take disjoint open neighborhoods
$U,V$ of $p,q$. Derivatives commute with restriction to an open subspace:
$U^{(j)}=U\cap X^{(j)}$, and similarly for $V$. In particular both $U$ and $V$
have last derivative rank $h$.

On $U$, choose the parity target with bottom color $0$. On $V$, choose the parity
target with bottom color $1$. Define the target arbitrarily outside $U\cup V$.
By \cref{lem:parity}, orientation $0$ requires $h+2$ layers on $V$, while
orientation $1$ requires $h+2$ layers on $U$. By
\cref{lem:profile-restrict}, both global orientations require at least $h+2$
layers. The upper bound already proved gives equality. Applying
\cref{thm:normal} finishes the proof.
\end{proof}

\begin{remark}[What information survives]
At finite height, the invariant sees the last derivative rank and whether the
terminal derivative has one point or more than one. It sees neither the exact
larger cardinality of that derivative nor the cardinalities of lower levels.
Compactness and metrizability are not needed for \cref{thm:height}.
\end{remark}

\section{Exact values on ordinal intervals}\label{sec:ordinals}

For $h,m\geq1$ finite, put
\[
 X_{h,m}=[0,\omega^h\cdot m].
\]
These are countable compact ordinal spaces. Standard background on ordinal
spaces and Cantor--Bendixson characteristics is discussed in \cite{AM}; the
specific derivative calculation needed here is elementary.

\begin{lemma}\label{lem:ordinal-rank}
For $1\leq j\leq h$,
\begin{equation}\label{eq:ordinal-derivative}
 X_{h,m}^{(j)}=
 \{\omega^j\eta:0<\eta\leq\omega^{h-j}\cdot m\}.
\end{equation}
In particular,
\[
 X_{h,m}^{(h)}=\{\omega^h,\omega^h\cdot2,\ldots,\omega^h\cdot m\},
 \qquad X_{h,m}^{(h+1)}=\varnothing.
\]
\end{lemma}
\begin{proof}
The nonisolated points of an ordinal interval are exactly its nonzero limit
ordinals. Below $\omega^h\cdot m$, these are precisely the positive multiples
of $\omega$. This gives the formula for $j=1$.

Within the ordered set of positive multiples of $\omega^j$, a point with
successor coefficient is isolated, including the first coefficient. A point
with limit coefficient is a limit of smaller such multiples. The nonzero limit
coefficients are the multiples of $\omega$. Therefore the next derivative
consists exactly of the positive multiples of $\omega^{j+1}$ in the indicated
range. This proves the induction. At step $h$ the remaining finite set is
discrete, so the next derivative is empty.
\end{proof}

\begin{corollary}[Finite-exponent ordinal formula]\label{cor:ordinal}
For positive finite $h,m$,
\begin{equation}\label{eq:ordinal-sm}
 \sm([0,\omega^h\cdot m])=
 \begin{cases}
  \ceil{\log_2(h+1)},&m=1,\\
  \ceil{\log_2(h+2)},&m\geq2.
 \end{cases}
\end{equation}
\end{corollary}

\begin{center}
\begin{tabular}{@{}rrrr@{}}
\toprule
Last rank $h$ & One terminal point & At least two & Duplication increases $\sm$\\
\midrule
0 & 0 & 1 & yes\\
1 & 1 & 2 & yes\\
2 & 2 & 2 & no\\
3 & 2 & 3 & yes\\
4 & 3 & 3 & no\\
5 & 3 & 3 & no\\
6 & 3 & 3 & no\\
7 & 3 & 4 & yes\\
8 & 4 & 4 & no\\
15 & 4 & 5 & yes\\
31 & 5 & 6 & yes\\
63 & 6 & 7 & yes\\
\bottomrule
\end{tabular}
\end{center}
The row $h=0$ refers to arbitrary discrete spaces, not the notation
$[0,\omega^0\cdot m]$. A singleton has value $0$; any larger discrete space has
value $1$.

\paragraph{Explicit hard targets.}
On $X_{h,1}$, color a point according to the parity of its Cantor--Bendixson
rank. In Cantor normal form, the rank of a nonzero point is its last nonzero
exponent; the point $0$ has rank $0$. On two copies, reverse all colors in one
copy. The canonical closure lengths are then $h+2$ in both orientations.
This is an explicit witness, not merely an existence argument using the
supremum over targets.

\paragraph{Why two copies form another ordinal space.}
For $h\geq1$, the intervals
\[
 [0,\omega^h],\qquad (\omega^h,\omega^h\cdot2]
\]
are clopen in $[0,\omega^h\cdot2]$ and are each homeomorphic to
$[0,\omega^h]$. The second interval has the same order type because removing
the first point from an infinite indecomposable ordinal $\omega^h$ does not
change that order type. Thus
\[
 X_{h,1}\sqcup X_{h,1}\cong X_{h,2}.
\]

\section{Sharp failures at every finite level, and sums more generally}

\begin{theorem}[Sharpness at every positive finite level]\label{thm:sharp}
For every integer $k\geq1$, let
\[
 K_k=[0,\omega^{\,2^k-1}].
\]
Then
\begin{equation}\label{eq:sharp}
 \sm(K_k)=k,\qquad
 \sm\!\left(\bigsqsum_{i<\kappa}K_k\right)=k+1
 \quad\text{for every cardinal }\kappa\geq2.
\end{equation}
Thus the one-extra-move upper bound for sums is attained at every positive
finite level, including by sums of two compact metrizable ordinal spaces.
\end{theorem}
\begin{proof}
Put $h=2^k-1$. Derivatives of a disjoint union are computed componentwise.
Each component has last rank $h$ and one terminal point; a sum of at least two
has the same last rank and at least two terminal points. By
\cref{thm:height}, the respective values are
\[
 \ceil{\log_2(2^k)}=k,\qquad
 \ceil{\log_2(2^k+1)}=k+1.
\]
\end{proof}

\begin{corollary}[Exact duplication threshold]\label{cor:threshold}
If $X$ is Hausdorff of last finite derivative rank $h$ and
$|X^{(h)}|=1$, then
\[
 \sm(X\sqcup X)>
 \sm(X)
 \quad\Longleftrightarrow\quad h+1\text{ is a power of }2.
\]
In that case the increase is exactly one.
\end{corollary}
\begin{proof}
The comparison is between $\ceil{\log_2(h+1)}$ and
$\ceil{\log_2(h+2)}$. Consecutive positive integers cross a ceiling-logarithm
threshold exactly when the smaller is a power of two.
\end{proof}

\subsection{Bounded finite heights}

Let $\{X_i:i\in I\}$ be any nonempty family of nonempty Hausdorff spaces with
finite last derivative ranks $h_i$. Suppose $H=\sup_i h_i<\omega$. Then $H$
is attained and
\[
 \left(\bigsqsum_iX_i\right)^{(H)}
 =\bigsqsum_{h_i=H}X_i^{(H)}.
\]
If there is a unique component of height $H+1$ and that component has a
singleton terminal derivative, the sum has value $\ceil{\log_2(H+1)}$.
In all other cases its value is $\ceil{\log_2(H+2)}$.
These two formulae can give the same numerical value.

\subsection{Unbounded finite heights}

The next upper bound is useful because the number of components need not be
countable.

\begin{proposition}[Parallel play with finite-time certification]\label{prop:omega-sum}
If every component $X_i$ has finite $\sm(X_i)$, then
\[
 \sm\!\left(\bigsqsum_iX_i\right)\leq\omega.
\]
If the component values are unbounded among the integers, equality holds.
\end{proposition}
\begin{proof}
Fix a target $A$ in the sum and, for each $i$, a finite winning strategy for
$A_i=A\cap X_i$. We describe countably many rounds, each using three questions.

In round $n$, first ask the union of the next local questions in every component
whose local strategy has not finished. Use a dummy open question in components
that have finished. The single global answer advances all unfinished local
simulations along the same answer history.

Next ask the union of the whole components whose current local simulation has
finished with output color $0$. Finally ask the analogous union for color $1$.
These are open sets because whole components are open. At either of these
certification questions, a yes answer leaves a monochromatic state, and the
Seeker has won. Empty local states may be assigned either output color.

Suppose all certification answers are no. Any point surviving the final
intersection would lie in a particular component $X_i$. Its local strategy
finishes after finitely many rounds. In that round or an earlier one, a
certification question includes the whole component $X_i$, and a no answer
removes the point. This is a contradiction. Thus the final state is empty.

The three questions per round have total order type $3\cdot\omega=\omega$.
For the lower bound when local values are unbounded, use subspace monotonicity.
\end{proof}

In particular, for any family of finite-height Hausdorff spaces whose last
ranks are unbounded in $\omega$, the sum has membership number exactly $\omega$.
This conclusion is independent of the cardinality of the index set.

\section{Finite products and the information lost by \texorpdfstring{$\sm$}{sm}}

\begin{lemma}[Product derivatives]\label{lem:product}
Let $X,Y$ be Hausdorff spaces of finite last derivative ranks $h,k$.
For finite $n$,
\begin{equation}\label{eq:product-derivative}
 (X\times Y)^{(n)}=
 \bigcup_{i+j=n}X^{(i)}\times Y^{(j)}.
\end{equation}
Thus the last rank of the product is $h+k$, and its terminal derivative is
$X^{(h)}\times Y^{(k)}$.
\end{lemma}
\begin{proof}
For closed subspaces $F\subseteq X$, $G\subseteq Y$, a point of $F\times G$
is isolated precisely when both coordinates are isolated in their respective
subspaces. Therefore
\[
 (F\times G)'=(F'\times G)\cup(F\times G').
\]
Derivation distributes over finite unions of closed sets: if a point is
isolated or absent in every one of finitely many closed sets, intersect finitely
many suitable neighborhoods to witness that it is not an accumulation point
of their union. The reverse inclusion is immediate. Induction proves
\eqref{eq:product-derivative}. At $n=h+k$ only the term with $i=h,j=k$ can be
nonempty, and all terms vanish at the next stage.
\end{proof}

\begin{corollary}[Exact product formula]\label{cor:product}
For nonempty finite-height Hausdorff spaces $X_1,\ldots,X_r$, put
$H=\sum_j h_j$. Then
\[
 \sm\!\left(\prod_{j=1}^rX_j\right)=
 \begin{cases}
 \ceil{\log_2(H+1)},&\text{every }X_j^{(h_j)}\text{ is a singleton},\\
 \ceil{\log_2(H+2)},&\text{otherwise}.
 \end{cases}
\]
In particular, for finite $r\geq1$,
\begin{equation}\label{eq:power}
 \sm\bigl([0,\omega]^r\bigr)=\ceil{\log_2(r+1)}.
\end{equation}
\end{corollary}

Knowing only the values of $\sm$ is not sufficient to compute product values.
For example, let
\[
 P=[0,\omega^2],\qquad Q=[0,\omega^3],\qquad S=[0,\omega].
\]
Then $\sm(P)=\sm(Q)=2$, but
\[
 \sm(P\times S)=2,\qquad \sm(Q\times S)=3.
\]
The finer data consisting of the last rank and singleton-versus-multiple
terminal derivative are compositional for finite products and bounded-height
sums. The logarithmic compression to $\sm$ discards precisely the numerical
information needed in this example.

\section{Countable spaces and the transfinite-height boundary}
\label{sec:countable}

The finite formula is not a transfinite ordinal-logarithm formula. On countable
scattered spaces, all infinite Cantor--Bendixson heights lead to the same game
value, namely $\omega$.

The point-separating number $\ps(X)$ is defined by the same game, but with
winning condition $|W_\beta|\leq1$ instead of monochromaticity for a target
\cite{CCS}.

\begin{lemma}[A countable $T_1$ upper bound]\label{lem:countable}
If $X$ is countable and $T_1$, then $\sm(X)\leq\omega$. If $X$ is also infinite,
then $\ps(X)=\omega$.
\end{lemma}
\begin{proof}
Enumerate $X=\{x_0,x_1,\ldots\}$. Ask the open complements
$X\setminus\{x_n\}$ successively. A no answer leaves at most the singleton
$\{x_n\}$. If every answer is yes, the final state is empty. This proves both
upper bounds. A finite binary tree has finitely many leaves and cannot separate
all points of an infinite space, proving the point-separation lower bound.
\end{proof}

\begin{theorem}[Countable scattered classification]\label{thm:countable}
Let $X$ be a nonempty countable scattered Hausdorff space. If its
Cantor--Bendixson height is finite, \cref{thm:height} applies. Otherwise
\begin{equation}\label{eq:infiniteheight}
 \sm(X)=\omega.
\end{equation}
\end{theorem}
\begin{proof}
Only the lower bound in the infinite-height case remains. For every finite $h$,
$X^{(h)}$ is nonempty and scattered, hence has an isolated point $p$.
Choose an open neighborhood $U$ in $X$ with $U\cap X^{(h)}=\{p\}$.
Derivatives commute with open restriction, so
\[
 U^{(h)}=\{p\},\qquad U^{(h+1)}=\varnothing.
\]
By \cref{thm:height} and subspace monotonicity,
$\sm(X)\geq\ceil{\log_2(h+1)}$. These integers are unbounded as $h$ varies.
Together with \cref{lem:countable}, this yields \eqref{eq:infiniteheight}.
\end{proof}

\begin{corollary}[All countable exponent intervals]\label{cor:allordinals}
Let $0<\alpha<\omega_1$ and $1\leq m<\omega$. Then
\[
 \sm([0,\omega^\alpha\cdot m])=
 \begin{cases}
 \ceil{\log_2(h+1)},&\alpha=h<\omega,\ m=1,\\
 \ceil{\log_2(h+2)},&\alpha=h<\omega,\ m\geq2,\\
 \omega,&\alpha\geq\omega.
 \end{cases}
\]
Every one of these spaces has $\ps=\omega$.
\end{corollary}
\begin{proof}
The finite-exponent cases were proved in \cref{cor:ordinal}. When
$\alpha\geq\omega$, the interval contains $[0,\omega^h]$ as a clopen subspace
for every positive finite $h$. Their membership numbers are unbounded, while
the whole space is countable and $T_1$. Use \cref{lem:countable}.
\end{proof}

For example,
\[
 \sm([0,\omega^\omega])
 =\sm([0,\omega^{\omega^2}])
 =\sm([0,\omega^{\omega^\omega}])=\omega.
\]
Thus extremely different countable order types can have the same membership
number. This is a mathematical limitation of the invariant, not a limitation
of the ordinal notation.

The countable scattered classification also covers every countable compact
Hausdorff space. For completeness, the fact that such a space is scattered is
proved in \cref{app:compact}, rather than invoked as an unproved classification
of countable compact spaces by ordinals.

\section{A purely order-theoretic form for finite posets}\label{sec:posets}

Equip a finite poset $P$ with the Alexandrov topology of upper sets. This is
$T_0$, and it is Hausdorff only when the order is an antichain. It is important
not to mistake this topology for the order topology on the infinite ordinal
spaces studied above.

Fix a target $A\subseteq P$ and color points by their membership in $A$.
A strictly increasing chain is \emph{alternating} if consecutive colors differ.
Let $a_b$ be the largest number of elements in an alternating chain whose
greatest element has color $b$, with $a_b=0$ if no point has that color.

\begin{theorem}[Alternating-chain formula]\label{thm:posets}
For every nonempty finite poset $P$ and target $A$,
\begin{equation}\label{eq:poset-profile}
 N_0(A,P)=1+a_1,\qquad N_1(A,P)=1+a_0,
\end{equation}
and therefore
\begin{equation}\label{eq:poset-sm}
 \boxed{\ \sm(A,P)=\ceil{\log_2(1+\min(a_0,a_1))}.\ }
\end{equation}
\end{theorem}
\begin{proof}
Closure in the upper-set topology is downward closure. We claim that
$F^c_n\ne\varnothing$ exactly when there is an alternating chain of length
$n$ ending in color $1-c$.

More precisely, $F^c_n$ is the downward closure of the least elements of all
such chains. This is immediate for $n=1$. In the next step, intersect with color
$(c+n+1)\bmod2$ and take downward closure. A point of that color below the
least element of an existing length-$n$ alternating chain has the opposite
color from that least element, so the inequality is strict and it can be
prepended to the chain. Conversely, the first point of any length-$(n+1)$
chain is produced this way. The induction follows.

The largest index of a nonempty $F^c_n$ is consequently $a_{1-c}$, giving
\eqref{eq:poset-profile}. Now apply \cref{thm:normal,thm:greedy}.
\end{proof}

\begin{corollary}[Quadratic-time computation]
Given the comparability relation of an $n$-element poset and a target, its
optimal finite query depth is computable in $O(n^2)$ time after choosing a
linear extension.
\end{corollary}
\begin{proof}
Process points in a linear extension. Set
\[
 a(x)=1+\max\{a(y):y<x,\ y\text{ has the opposite color to }x\},
\]
where the maximum of the empty set is $0$. Then
$a_b=\max\{a(x):x\text{ has color }b\}$. Each comparable pair is inspected at
most once. Formula \eqref{eq:poset-sm} completes the computation.
\end{proof}

\begin{example}[A finite orientation conflict]
An alternating four-element chain has profile $(4,5)$ or $(5,4)$, depending on
its orientation. Its target takes two questions. Two disjoint chains with
opposite orientations have profile $(5,5)$, so their combined target takes
three questions. The archive contains an explicit optimal query tree for this
eight-point example.
\end{example}

\begin{remark}
This order-theoretic formula is not obtained by applying the Hausdorff
finite-height theorem to a non-Hausdorff space. Both results are separate
consequences of the general finite normal form, and their hypotheses are kept
distinct.
\end{remark}

\section{Algorithms, exact checks, and reproducibility}\label{sec:checks}

The proof of \cref{thm:normal} gives a strategy compiler. First construct the two
canonical closure sequences; choose the shorter one; then binary-search its
layer indices. At a node where the possible indices are $a,\ldots,b-1$, choose
$j=\lfloor(a+b)/2\rfloor$ and ask $X\setminus F_j$. A yes answer selects
$a,\ldots,j-1$, and a no answer selects $j,\ldots,b-1$.

For a general infinite space, computing a closure is a mathematical operation,
not necessarily an effective algorithm on a presentation of the space. The
complexity claim in \cref{sec:posets} is specifically for finite posets whose
comparability relation is supplied.

\subsection{Independent verification paths}

The program \texttt{code/verify.py} uses only the Python standard library.
For each finite poset and every target, it compares three independently
implemented calculations: exhaustive minimax over all upper-set questions,
canonical downward-closure profiles, and the alternating-chain dynamic program.
It also constructs a balanced strategy tree from the shorter closure chain and
checks the color returned on every point.

The exhaustive minimax recurrence is
\[
 d_A(S)=
 \begin{cases}
 0,&S\text{ is monochromatic},\\
 1+\displaystyle\min_U\max\{d_A(S\cap U),d_A(S\setminus U)\},&\text{otherwise},
 \end{cases}
\]
where $U$ ranges over upper sets splitting $S$ nontrivially. States are encoded
as bit masks, and the recursion is memoized. This recurrence does not use the
closure formula being tested.

\begin{center}
\begin{tabular}{@{}rrr@{}}
\toprule
Number of points & Naturally labelled posets & Targets checked\\
\midrule
1 & 1 & 2\\
2 & 2 & 8\\
3 & 7 & 56\\
4 & 40 & 640\\
5 & 357 & 11,424\\
6 & 4,824 & 308,736\\
\midrule
Total & 5,231 & 320,866\\
\bottomrule
\end{tabular}
\end{center}

All checks passed. ``Naturally labelled'' means that every strict comparison
$i<_{P}j$ satisfies $i<j$ as an integer. These are not counts of isomorphism
classes. Every finite poset has at least one such labelling, so the search
covers every isomorphism type through six points, with repetitions possible.

The included run took approximately 13 seconds in the execution environment
used for this report. That observation is not a portable performance guarantee.
The machine-readable result is in \texttt{data/verification.json}; the text log,
count table, ordinal-value table, sharp-family table, and eight-point strategy
certificate are provided alongside it.

\subsection{Scope of the computational evidence}

Finite upper-set topologies are not convergent sequences and do not model the
Hausdorff accumulation behavior of ordinal intervals. Accordingly, the finite
checks do \emph{not} establish the infinite-space counterexample, the
Cantor--Bendixson formula, or its transfinite consequences. Those results follow
from the proofs in \cref{sec:small,sec:height,sec:countable}. The computation
instead checks the general finite game/closure/alternation mechanism in a
large family where a direct independent minimax computation is possible.

To reproduce the included exhaustive run, execute
\begin{lstlisting}
python3 code/verify.py --max-n 6
\end{lstlisting}
To rebuild the PDF, execute \texttt{make pdf}, or run \texttt{pdflatex} three times on
\texttt{article.tex}. All inputs needed for compilation are local; no network
access, external bibliography processor, or data regeneration is required.

\section{What is settled and what is not}

The topological-sum identity \eqref{eq:question} is false as stated. The
counterexample can be expressed using the ordinary ordinal intervals
$[0,\omega]$ and $[0,\omega\cdot2]$, and the one-move defect is sharp at every
positive finite level. The proof also yields exact computations for every
Hausdorff space of finite Cantor--Bendixson height and every countable scattered
Hausdorff space.

The two central distinctions are structural. First, the number of alternating
layers is different from the number of adaptive questions: they are related
by a ceiling logarithm. Second, the minimum over the two orientations is
different from the pair of oriented lengths. Taking a disjoint union can force
a common orientation and increase the minimum length by one; when that passes
a power of two, it costs a full additional query.

No transfinite version of \eqref{eq:logdepth} is asserted. In particular, the
symbol $\infty$ in the finite profile is not an ordinal rank and cannot simply
be inserted into an ordinal logarithm. Nothing here determines the games at
arbitrary uncountable lengths or resolves the other problems posed in
\cite{CCS}. Those are separate questions.

The research claim should be read at its demonstrated scale: an explicit
negative solution to an explicitly stated question, together with a complete additional
classification on a natural order-topological class. The proofs are intended
to be independently checkable; their priority and presentation can be assessed
separately from their mathematical content.

\appendix
\section{Why countable compact Hausdorff spaces are scattered}\label{app:compact}

This appendix supplies the fact used after \cref{cor:allordinals} without
requiring the full ordinal classification of countable compact spaces.

\begin{lemma}
A nonempty compact Hausdorff space with no isolated points has at least
$2^{\aleph_0}$ points.
\end{lemma}
\begin{proof}
Construct nonempty open sets $U_s$ indexed by finite binary strings. Start with
the whole space. Inside each $U_s$, select two distinct points; this is possible
because no nonempty open set is a singleton. By regularity, choose nonempty
open neighborhoods $U_{s0},U_{s1}$ with disjoint closures contained in $U_s$.

For each infinite binary string $z$, the compact sets
$\overline{U_{z\upharpoonright n}}$ are nonempty and nested, so their intersection
is nonempty. Choose one point from each intersection. Distinct binary strings
split at some level and hence give points in disjoint closed sets. This produces
an injection of $2^\N$ into the space.
\end{proof}

\begin{proposition}
Every countable compact Hausdorff space is scattered.
\end{proposition}
\begin{proof}
If a space is not scattered, some nonempty subspace $B$ has no isolated points.
Its closure $F$ has no isolated points either: an isolated point of $F$ would,
by density, belong to $B$ and be isolated in $B$. In a compact Hausdorff space,
$F$ is compact Hausdorff. The preceding lemma would make $F$ uncountable, a
contradiction.
\end{proof}

This proof also shows why no additional set-theoretic hypothesis such as the
continuum hypothesis is needed. Uncountability of $2^\N$ is sufficient.

\section{A proof audit and implementation checklist}

The following points isolate possible sources of off-by-one or topology errors.

\paragraph{Zero questions.}
A constant target has $L=1$, not $L=0$, on a nonempty space. Thus its depth is
$\ceil{\log_2 1}=0$. The empty space is assigned $\sm=0$ separately.

\paragraph{Derivative rank.}
The integer $h$ is the index of the \emph{last nonempty} derivative, not the
height. The height is $h+1$. A convergent sequence has $h=1$ and height $2$.

\paragraph{First closure.}
The general estimate is $F_n\subseteq X^{(n-1)}$ for $n\geq1$, not
$F_n\subseteq X^{(n)}$. The extra first step is exactly what creates the two
oriented lengths $h+1$ and $h+2$ for a parity target.

\paragraph{Hausdorff separation.}
The lower bound with two terminal points uses disjoint open neighborhoods.
The finite-height theorem is therefore stated for Hausdorff spaces. Its proof
is not silently applied to arbitrary $T_0$ posets.

\paragraph{Arbitrary answers.}
The binary-search proof wins on every answer sequence, including one whose
state becomes empty. In the $\omega$-length certification strategy, an
all-no certification path has empty final intersection. There is no hidden
assumption that the Hider fixes a point beforehand.

\paragraph{Infinite sums.}
In the oriented finite-profile formula, an unbounded supremum of integers is
recorded as $\infty$, meaning ``no finite decomposition''. The separate
certification argument is needed to obtain the actual ordinal value $\omega$.

\paragraph{Independent checking.}
The minimax implementation enumerates all open queries. It does not import the
closed-layer answer as its return value. The mathematical proof and the
finite computational check are deliberately distinct.

\begin{thebibliography}{9}
\small
\setlength{\itemsep}{0.35em}
\setlength{\parskip}{0pt}

\bibitem{CCS}
Lucas Chiozini, Tam\'as Csern\'ak, and Lajos Soukup,
\emph{Gamification of the $T_0$-pseudoweight via cut-and-choose games on
topological spaces}, arXiv:2510.05754v3, 29 May 2026.
The arXiv abstract record retains the shorter title
\emph{Cut-and-choose games in topological spaces}.
\url{https://arxiv.org/abs/2510.05754v3}.
Problem 1.3 and Theorem 3.4 are the principal source references.
Accessed 19 September 2026.

\bibitem{BGKS}
C\'elia Borlido, Mai Gehrke, Andreas Krebs, and Howard Straubing,
\emph{Difference hierarchies and duality with an application to formal languages},
Topology and its Applications \textbf{273} (2020), 106975.
\url{https://doi.org/10.1016/j.topol.2019.106975}.
Preprint: \url{https://arxiv.org/abs/1812.01921}.
Cited for the classical difference-chain framework, not as a source for the
new game computations in this report.

\bibitem{AM}
Borys \'Alvarez-Samaniego and Andr\'es Merino,
\emph{Some properties related to the Cantor--Bendixson derivative on a Polish space},
New Zealand Journal of Mathematics \textbf{50} (2020), 207--218.
\url{https://doi.org/10.53733/82}.
Preprint: \url{https://arxiv.org/abs/2003.01512}.
The volume PDF is dated 2020; the journal landing page records online
publication in 2021. Cited for ordinal-space and derivative background.

\end{thebibliography}
\end{document}
